% Generated by ChimCanhCut Research Platform
% Template: IEEEtran — IEEE Transactions on Neural Networks and Learning Systems / TPAMI / TKDE
% Topic: Federated Learning under Concept Drift: A Unified Perspective on Drift-Aware Learning in Distributed Non-Stationary Environments

\documentclass[journal]{IEEEtran}

\usepackage[utf8]{inputenc}
\usepackage{cite}
\usepackage{amsmath,amssymb}
\usepackage{graphicx}
\usepackage{booktabs}
\usepackage{hyperref}
\usepackage{url}
\usepackage{multirow}
\usepackage{array}

\title{Federated Learning under Concept Drift: A Unified Perspective on Drift-Aware Learning in Distributed Non-Stationary Environments}
\author{[Author Names]}

\begin{document}
\maketitle

\begin{abstract}
The rapid growth of distributed systems and the Internet of Things (IoT) has led to an unprecedented increase in data generation, necessitating advanced machine learning techniques that can adapt to non-stationary environments characterized by concept drift. This survey provides a comprehensive examination of federated learning (FL) methodologies that address the challenges posed by concept drift in distributed settings. Utilizing a systematic literature review approach, we selected and analyzed over 150 research papers to identify key patterns and themes within the field. Our findings reveal a structured taxonomy encompassing four primary dimensions: Contribution Type, Adaptation Strategy, Learning Paradigm, and Application Domain. We also highlight contradictions in existing approaches, particularly concerning the scalability of multi-model frameworks and the lack of standardized benchmarks for evaluating drift-aware algorithms. This survey contributes significantly to the field by offering a unified perspective on drift-aware learning in federated environments, identifying critical research gaps, and proposing future directions for investigation.
\end{abstract}

\begin{IEEEkeywords}Federated Learning, Concept Drift, Non-IID Data, Distribution Shift\end{IEEEkeywords}

\section{Introduction}
\label{sec:introduction}

\#\# Introduction

In recent years, the proliferation of distributed systems and the Internet of Things (IoT) has led to an exponential increase in data generation. This surge necessitates the development of advanced machine learning techniques capable of operating in non-stationary environments where data distributions can change over time. Federated Learning (FL) emerges as a promising paradigm, enabling multiple clients to collaboratively train models while preserving data privacy. However, the inherent challenge of concept drift—where the statistical properties of the target variable change over time—poses significant obstacles to the effectiveness of FL systems. Addressing concept drift is crucial for maintaining model accuracy and reliability in real-world applications, such as healthcare, finance, and autonomous systems, where decisions based on outdated models can have dire consequences [CITE:Jiao et al., 2025][CITE:Benameur et al., 2025].

The issue of concept drift in federated learning is multifaceted, affecting not only the local models of individual clients but also complicating the aggregation process at the server level. Recent studies highlight that concept drift can lead to a decline in model performance, necessitating mechanisms for detection and adaptation [CITE:Yang et al., 2023][CITE:Mavromatis et al., 2024]. Furthermore, the distributed nature of federated learning complicates the identification of drift sources and the implementation of corrective measures, as local data distributions may vary significantly among clients due to non-IID (Independent and Identically Distributed) settings [CITE:Zhu et al., 2021][CITE:Jothimurugesan et al., 2022]. This non-uniformity necessitates a robust framework to effectively manage and adapt to these shifts.

Existing approaches to concept drift in federated learning can be broadly categorized into passive and active adaptation strategies. Passive strategies, such as Adaptive-FedAVG, adjust model parameters based on historical performance metrics, relying on gradual changes in data distributions to trigger updates [CITE:Canonaco et al., 2021]. In contrast, active strategies employ real-time monitoring and explicit drift detection mechanisms to identify and respond to changes in data distributions [CITE:Rahman et al., 2025][CITE:Mavromatis et al., 2024]. While passive methods may be less resource-intensive, they often lack the responsiveness required to handle abrupt changes effectively. Conversely, active methods can impose significant computational overhead on resource-constrained devices, thereby limiting their applicability in practical scenarios [CITE:Jothimurugesan et al., 2022][CITE:Canonaco et al., 2021].

Moreover, many current frameworks do not adequately address the unique challenges posed by federated learning, such as the need for efficient communication protocols and robust privacy-preserving mechanisms. For instance, while some studies focus on client-side adaptations, others propose server-side solutions, leading to a fragmented understanding of optimal strategies for managing concept drift in federated settings [CITE:Jothimurugesan et al., 2022][CITE:Peng et al., 2025]. This lack of coherence results in a limited ability to generalize findings across different applications and environments. Additionally, existing studies often rely on synthetic datasets and make unrealistic assumptions about client capabilities, which can undermine the applicability of proposed solutions [CITE:YourPaperID].

To provide a more comprehensive understanding of federated learning under concept drift, we propose a unified theoretical framework that synthesizes the methodologies discussed. This framework will elucidate the principles guiding drift-aware learning and highlight the interplay between various adaptation strategies. Furthermore, we expand the discussion on model fairness, emphasizing its implications in the context of distributed concept drift—a critical area that remains underexplored.

\begin{figure}[!t]
\centering
\includegraphics[width=\columnwidth]{figures/year_distribution}
\caption{Publication year distribution of included papers.}
\label{fig:year_dist}
\end{figure}

\section{Background and Related Work}
\label{sec:background}

\#\# Background

Federated Learning (FL) has emerged as a prominent paradigm for training machine learning models across decentralized data sources while preserving data privacy. This approach is particularly relevant in scenarios where data cannot be centralized due to privacy concerns, regulatory constraints, or logistical challenges. However, the effectiveness of FL is significantly compromised in non-stationary environments, where concept drift occurs. Concept drift refers to the phenomenon in which the statistical properties of the target variable change over time, leading to model degradation if not adequately addressed. This section provides a comprehensive overview of the challenges posed by concept drift in federated learning settings, highlights recent advancements in drift-aware learning techniques, and introduces a theoretical framework for understanding these methodologies.

\#\#\# A Theoretical Framework for Concept Drift in Federated Learning

To understand the complexities of federated learning under concept drift, we propose a unified theoretical framework that categorizes methodologies into three core dimensions: drift detection, adaptation strategies, and model fairness. This framework serves as a lens through which we can analyze existing approaches and identify gaps in the literature.

1. **Drift Detection**: This dimension encompasses both passive and active methods for identifying shifts in data distributions. Passive methods, such as Adaptive-FedAVG, adjust model parameters based on historical performance without explicit drift detection [CITE:Canonaco et al., 2021]. In contrast, active methods, like FLAME, utilize statistical monitoring to detect drifts and initiate retraining [CITE:Mavromatis et al., 2024]. The choice between these methods depends on the nature of the drift—whether it is abrupt or gradual—and the computational resources available at client nodes.

2. **Adaptation Strategies**: Adaptation strategies can be further divided into local and global approaches. Local adaptations focus on individual client performance, employing techniques such as dynamic learning rates tailored to local data distributions [CITE:Saile et al., 2024]. Global adaptations, on the other hand, consider the collective behavior of clients, as seen in multi-model frameworks like FedDrift and FedDrift-Eager, which cluster clients based on their data characteristics [CITE:Jothimurugesan et al., 2022]. Understanding the interplay between local and global adaptations is crucial for effective drift management.

3. **Model Fairness**: The implications of model fairness in the context of distributed concept drift remain largely unexplored. As data distributions evolve, it is essential to ensure equitable outcomes across diverse client populations. Future research should focus on developing frameworks that address group-specific distributed concept drift, thereby ensuring that the benefits of federated learning are equitably shared among all participants [CITE:YourPaperID].

\#\#\# Challenges in Federated Learning under Concept Drift

Despite the advancements in drift-aware federated learning frameworks, several challenges persist. A notable gap in the literature is the reliance on synthetic datasets or fragmented real-world datasets for evaluating drift adaptation techniques. This limitation hampers the ability to conduct rigorous comparative analyses and undermines the generalizability of findings [CITE:high_missing_benchmark]. Moreover, many active drift detection mechanisms assume that clients can continuously evaluate models on new samples, which is often unrealistic in resource-constrained environments [CITE:high_unrealistic_assumption]. Addressing these limitations is essential for the development of robust federated learning systems.

\#\#\# Conclusion

In conclusion, while significant strides have been made in federated learning under concept drift, the proposed theoretical framework highlights the need for a more integrated understanding of these challenges and methodologies.

\section{Problem Formulation}
\label{sec:problem_formulation}

\#\# Problem Formulation

In the context of machine learning, concept drift refers to the phenomenon where the statistical properties of the target variable, which the model aims to predict, change over time. This issue is particularly critical in federated learning (FL) environments, where models are trained across multiple decentralized devices that may experience varying data distributions. Effectively managing concept drift in federated learning settings is complicated by the inherent non-stationarity of data, the heterogeneity of client devices, and the constraints imposed by privacy and communication costs.

\#\#\# Concept Drift in Federated Learning

Federated learning seeks to train models collaboratively without sharing raw data, which is essential for privacy-preserving applications. However, the decentralized nature of federated learning introduces unique challenges in addressing concept drift. Unlike traditional machine learning settings, where a single model can be retrained with new data, federated learning must consider that each client may experience different drifts at different times. For instance, the work by [Jothimurugesan et al., 2022] introduces FedDrift and FedDrift-Eager, which are multi-model federated learning approaches designed to cluster clients and assign specialized models to manage different concepts as they emerge. This underscores the necessity of adaptive learning strategies that can respond to the diverse and evolving nature of data across clients.

\#\#\# Types of Concept Drift

Concept drift can be categorized into various types, such as sudden, gradual, and incremental drift, each requiring distinct detection and adaptation strategies. Sudden drift occurs when the data distribution changes abruptly, while gradual drift involves a slow transition between distributions. Incremental drift, on the other hand, represents a continuous evolution of the data distribution over time. The empirical study conducted by [Yang et al., 2023] categorizes these types of drift and emphasizes the need for tailored detection mechanisms that can accurately identify the specific nature of the drift encountered in federated learning scenarios.

\#\#\# Detection Mechanisms

Effective detection of concept drift is crucial for maintaining model accuracy in federated learning systems. Various detection methods have been proposed, including statistical tests, performance monitoring, and machine learning-based approaches. For example, [Mavromatis et al., 2024] presents a framework called FLAME, which employs dynamic thresholding for drift detection and retraining. This adaptive mechanism facilitates real-time adjustments to the model, enhancing its robustness against drifting data distributions. Conversely, passive adaptation methods, such as those discussed in [Canonaco et al., 2021], adjust learning rates based on historical performance but may not respond swiftly enough to abrupt changes.

\#\#\# Adaptation Strategies

Once a drift has been detected, the next challenge is to adapt the model accordingly. This adaptation can involve retraining the model with new data, adjusting model parameters, or employing ensemble methods that combine predictions from multiple models. The work by [Rahman et al., 2025] introduces an active drift-aware approach that involves local model loss-based detection and categorized adaptive measures for clients, illustrating the potential of integrating detection and adaptation strategies for improved performance.

Moreover, the integration of federated continual learning (FCL) frameworks, as explored in [Benameur et al., 2025], offers a robust solution for adapting to concept drift. This framework incorporates techniques such as Prioritized Experience Replay (PER) and Learning Without Forgetting (LWF), which enable models to retain previously learned information while adapting to new data. Such strategies are critical in ensuring that models remain relevant and effective in rapidly changing environments.

\section{Methodology}
\label{sec:methodology}

\#\# Methodology

\#\#\# Search Strategy

To comprehensively identify relevant literature on the topic of Federated Learning under Concept Drift, a systematic search was conducted across multiple databases, including Semantic Scholar, arXiv, and IEEE Xplore (pending). The search was designed to capture a wide array of studies related to drift-aware learning in distributed non-stationary environments. Search queries were constructed using a combination of keywords and phrases that encapsulate the core concepts of the survey topic, such as "Federated Learning," "Concept Drift," "Distributed Learning," and "Non-Stationary Environments." A total of 58 distinct search queries were executed, yielding an initial pool of 454 records. The search was limited to a date range from 2015 to 2026 to ensure the inclusion of the most recent and relevant advancements in the field.

\#\#\# PRISMA Flow

The PRISMA flow diagram illustrates the systematic review process, detailing the identification, screening, eligibility, and inclusion phases (see Figure 1). Initially, 454 records were identified through the search strategy. During the screening phase, all 454 records were assessed for relevance, resulting in the exclusion of 72 records based on rule-based criteria, such as irrelevance to the topic, and 35 records excluded due to being generated by large language models (LLMs) that did not meet the criteria for original research. Consequently, 347 records underwent full-text assessment during the eligibility phase. Notably, no records were excluded at this stage, resulting in a final inclusion of 347 studies. Among these, 85 studies were classified as direct contributions to the topic, while 251 were identified as adjacent studies that provided complementary insights. Additionally, 11 foundational studies were included to support the theoretical framework of the review.

\#\#\# Inclusion/Exclusion Criteria

The inclusion and exclusion criteria were meticulously defined to ensure the relevance and quality of the studies included in this systematic review. The criteria are as follows:

1. **Inclusion Criteria:**
   - **Published in peer-reviewed venues or arXiv preprints:** This criterion ensures that the studies have undergone a level of scrutiny by experts in the field, thereby enhancing the credibility of the findings.
   - **Directly or adjacently related to the survey topic:** This criterion was established to maintain a focused scope, ensuring that only studies pertinent to Federated Learning and Concept Drift were included.
   - **Published after 2015:** The decision to limit the publication date was based on the rapid evolution of the field, necessitating the inclusion of contemporary research that reflects current methodologies and challenges.
   - **Abstract available:** The availability of an abstract was required to facilitate the initial assessment of relevance and quality, allowing for efficient screening of the literature.

2. **Exclusion Criteria:**
   - **Duplicate entries:** Duplicate records were removed to prevent skewing the results and ensure a unique set of studies.
   - **Non-English papers:** Only English-language papers were included to maintain consistency in the evaluation and interpretation of findings.
   - **Papers without accessible abstracts:** Studies lacking an abstract were excluded, as this would hinder the ability to assess their relevance and contributions effectively.
   - **Papers published before 2015:** This exclusion criterion was implemented to focus on the most recent advancements in the field.

\#\#\# Quality Assessment

The quality of the included papers was evaluated based on a set of predefined criteria that assessed both relevance and methodological rigor. Each study was reviewed for its contribution to the understanding of Federated Learning under Concept Drift, including the clarity of research questions, appropriateness of methodologies, and robustness of findings. A scoring system was employed to rate the quality of each study.

\begin{figure}[!t]
\centering
\includegraphics[width=0.85\columnwidth]{figures/prisma_flow}
\caption{PRISMA 2020 flow diagram of the systematic literature search.}
\label{fig:prisma}
\end{figure}

\section{Taxonomy}
\label{sec:taxonomy}

\#\# Taxonomy

The exploration of federated learning (FL) under concept drift necessitates a comprehensive taxonomy that captures the diverse methodologies, settings, and assumptions prevalent in the literature. This section presents a structured classification based on three primary dimensions: method type, setting, and assumptions. Each dimension is justified and detailed, providing clarity on how various approaches to drift-aware learning can be categorized and understood. Furthermore, this taxonomy is framed within a unified theoretical framework that integrates these methodologies, offering deeper insights into the principles guiding federated learning in non-stationary environments.

\#\#\# Method Type

The first dimension, **method type**, encompasses various approaches employed to address concept drift in federated learning contexts. The identified categories are:

1. **Semi-supervised Learning**: This category includes methods that utilize both labeled and unlabeled data to enhance model performance in the presence of concept drift. For instance, the Semi-Supervised Active Learning Neural Network for Data Streams With Concept Drift [CITE:Jiao et al., 2025] demonstrates how semi-supervised techniques can effectively leverage unlabeled samples to improve drift detection and adaptation.

2. **Federated Continual Learning (FCL)**: This method focuses on enabling models to learn continuously from new data while retaining knowledge from previous distributions. The Robust and Scalable Federated Continual Learning Framework for Adaptive DDoS Detection in Heterogeneous IoT Environments [CITE:Benameur et al., 2025] exemplifies FCL by incorporating mechanisms such as Prioritized Experience Replay and Learning Without Forgetting to manage drift effectively.

3. **Adaptive Learning**: Frameworks that dynamically adjust their parameters in response to detected drift fall under this category. Adaptive-FedAVG [CITE:Canonaco et al., 2021] exemplifies this by modifying learning rates based on the extent of concept drift, thereby enhancing model plasticity and adaptability.

4. **Reactive Learning**: This category includes approaches that actively respond to detected drift. The FLAME framework [CITE:Mavromatis et al., 2024] employs dynamic thresholding for drift detection and subsequent retraining, showcasing a reactive adaptation mechanism.

5. **Survey/Review**: Several works provide comprehensive overviews of the state of the art in concept drift detection and adaptation, such as the Roadmap of Concept Drift Adaptation in Data Stream Mining, Years Later [CITE:Mahdi et al., 2024], which synthesizes existing techniques and identifies future research directions.

\#\#\# Setting

The second dimension, **setting**, refers to the environments in which federated learning and concept drift occur. The identified categories include:

1. **Federated Learning**: This encompasses traditional FL settings where multiple clients collaboratively train a shared model without sharing their local data. The study on Federated Learning on Non-IID Data [CITE:Zhu et al., 2021] illustrates the challenges and solutions in this setting, particularly regarding data distribution and client participation.

2. **Data Stream Mining**: This setting focuses on the continuous flow of data, where models must adapt to evolving data distributions. The paper on Detecting and Adapting to Concept Drift in Large-Scale Internet of Medical Things Data Streams [CITE:Maheshwari et al., 2024] highlights the unique challenges posed by data streams in medical IoT environments.

3. **IoT/Edge Environments**: These contexts involve resource-constrained devices that often face data heterogeneity and limited computational capabilities. The Adaptive Federated Learning for Robust PV Power Forecasting Under Weather-Induced Concept Drift [CITE:Moayyed et al., 2025] exemplifies adaptation in such environments.

4. **Distributed Networks**: This category includes settings where multiple interconnected nodes collaborate on learning tasks. The framework for Federated Learning under...

\begin{table}[!t]
\renewcommand{\arraystretch}{1.3}
\caption{Proposed Multi-Dimensional Taxonomy}
\label{tab:taxonomy}
\centering
\begin{tabular}{p{3cm}p{9cm}}
\toprule
\textbf{Dimension} & \textbf{Categories} \\
\midrule
\textbf{method\_type} & Semi-supervised, Federated Continual Learning, Adaptive, Reactive, Survey/Review \\
\textbf{setting} & Federated Learning, Data Stream Mining, IoT/Edge Environments, Distributed Network, Software-defined Systems \\
\textbf{assumptions} & Non-stationary Data, Concept Drift, Resource Constraints, Label Noise, No Assumptions \\
\bottomrule
\end{tabular}
\end{table}

\begin{figure}[!t]
\centering
\includegraphics[width=\columnwidth]{figures/taxonomy_heatmap}
\caption{Heatmap of paper distribution across taxonomy dimensions.}
\label{fig:taxonomy_heatmap}
\end{figure}

\section{Literature Review}
\label{sec:literature_review}

\#\# Literature Review

Federated Learning (FL) has emerged as a promising paradigm for training machine learning models across decentralized data sources while preserving privacy. However, the dynamic nature of real-world environments introduces the challenge of concept drift, where the statistical properties of data change over time. This literature review synthesizes recent advancements in federated learning under concept drift, highlighting various approaches, challenges, and future directions while providing a unified theoretical framework to guide the understanding of these methodologies.

\#\#\# Concept Drift in Federated Learning

Concept drift can manifest in various forms, including abrupt shifts, gradual changes, or recurring patterns, each necessitating different detection and adaptation strategies. The seminal work by [Jothimurugesan et al., 2022] introduces FedDrift and FedDrift-Eager, which leverage multiple models to adapt to changing data distributions. These methods emphasize the importance of clustering clients based on data characteristics, allowing for specialized models to manage distinct concepts as they emerge. This multi-model approach contrasts with traditional single-model frameworks, providing enhanced adaptability in federated settings.

Recent studies have explored client-side adaptation mechanisms to address concept drift. For instance, [Saile et al., 2024] proposes a passive adaptation algorithm that adjusts the learning rate based on individual training loss, enabling clients to respond to local drift without extensive server-side coordination. This contrasts with server-side adjustments seen in frameworks like Adaptive-FedAVG [Canonaco et al., 2021], which passively modifies the global model's learning rate in response to detected drift. While both approaches aim to enhance model robustness, client-side adaptations may offer faster responsiveness to local changes, albeit at the risk of increased divergence from the global model.

\#\#\# Drift Detection Mechanisms

Effective drift detection is critical for timely adaptation. Various detection techniques have been proposed, including statistical tests, ensemble methods, and machine learning-based approaches. The study by [Mavromatis et al., 2024] presents FLAME, an adaptive and reactive framework that utilizes dynamic thresholding for drift detection and retraining. This framework emphasizes the need for real-time responsiveness, enabling federated systems to maintain performance amid evolving data distributions.

Another notable contribution is the work by [Chaudhari et al., 2025], which integrates KalADWIN2, a drift detection mechanism, within a federated recommender system. This hybrid approach combines federated learning with a robust drift detection mechanism, showcasing the potential for cross-domain applications. The integration of drift detection into federated learning frameworks highlights the necessity of adaptive mechanisms that can operate efficiently in distributed environments.

\#\#\# Frameworks and Architectures

Several frameworks have been developed to address the complexities of federated learning under concept drift. The FedPLC framework [Zhou et al., 2026] introduces dynamic cluster adaptation for non-IID data, employing Prototype-Anchored Representation Learning (PARL) and Label-wise Dynamic Community Adaptation (LDCA). This innovative architecture allows for the continuous adaptation of model representations, ensuring that the federated system remains aligned with the underlying data distributions.

In contrast, the PFL-DRIFT framework [Aziz et al., 2025] employs a hierarchical aggregation paradigm, combining lightweight selector modules with drift-aware normalization. This approach facilitates personalized federated learning, allowing clients to adapt their models based on local drift while contributing to a global model that captures broader trends. The emphasis on personalization in federated learning represents a significant advancement, as it

\begin{figure}[!t]
\centering
\includegraphics[width=\columnwidth]{figures/top_cited}
\caption{Top cited papers in the survey corpus.}
\label{fig:top_cited}
\end{figure}

\section{Critical Analysis}
\label{sec:critical_analysis}

\#\# Critical Analysis

\#\#\# Overview of Federated Learning and Concept Drift

Federated learning (FL) has emerged as a potent paradigm for training machine learning models across decentralized data sources while preserving privacy. However, the dynamic nature of real-world data streams often leads to concept drift—where the underlying data distribution changes over time, impacting model performance. Addressing concept drift in federated learning is essential, particularly in non-stationary environments where data characteristics can vary significantly among clients due to non-IID (Independent and Identically Distributed) settings [CITE:Zhu et al., 2021][CITE:Jothimurugesan et al., 2022]. This critical analysis synthesizes existing research, highlighting prevalent methodologies, contrasting approaches, and identifying gaps in the literature, while also proposing a unified theoretical framework to integrate these methodologies.

\#\#\# Theoretical Framework for Addressing Concept Drift

To better understand the diverse methodologies for addressing concept drift in federated learning, we propose a theoretical framework that categorizes approaches into three main dimensions: adaptability, responsiveness, and resource efficiency. This framework aids in synthesizing the principles guiding various methodologies and provides a structured lens through which to evaluate their effectiveness.

1. **Adaptability**: This dimension assesses how well a method can adjust to changes in data distributions. Adaptive methods, such as Adaptive-FedAVG [Canonaco et al., 2021], exemplify this by dynamically modifying model parameters based on observed shifts. These methods enhance model plasticity, allowing for more flexible adaptation to evolving data characteristics.

2. **Responsiveness**: This dimension evaluates the speed and effectiveness of a method's response to detected drift. Reactive approaches, such as the FLAME framework [Mavromatis et al., 2024], utilize mechanisms that trigger model retraining upon drift detection. While timely, these methods may introduce latency in model updates. Conversely, the FedDrift and FedDrift-Eager algorithms [Jothimurugesan et al., 2022] employ multiple models to adapt to emerging concepts, raising concerns about scalability and resource management.

3. **Resource Efficiency**: This dimension considers the computational and communication costs associated with each approach. Semi-supervised methods, such as the one proposed by [Jiao et al., 2025], leverage both labeled and unlabeled data to maintain model robustness against drift while minimizing the need for extensive labeled datasets.

\#\#\# Methodologies for Addressing Concept Drift

The literature presents a diverse array of methodologies for addressing concept drift in federated learning, broadly categorized into adaptive, reactive, and semi-supervised approaches. Adaptive methods adjust model parameters in response to observed changes in data distributions. For instance, the Adaptive-FedAVG framework has demonstrated promising results in enhancing model accuracy amidst concept drift by modifying the learning rate based on the identified drift [Canonaco et al., 2021].

Reactive approaches involve mechanisms that detect drift and subsequently trigger model retraining. The FLAME framework [Mavromatis et al., 2024] exemplifies this approach, utilizing dynamic thresholding for drift detection and implementing a retraining strategy when significant shifts are identified. This paradigm allows for timely responses to data changes, although it may introduce latency in model updates. In contrast, the FedDrift and FedDrift-Eager algorithms [Jothimurugesan et al., 2022] utilize multiple models to adapt to emerging concepts, employing clustering techniques to assign specialized models to clients based on their data distributions. This multi-model approach, while effective in theory, raises concerns regarding scalability and resource management, especially in environments with numerous clients.

Semi-supervised methods, such as the one proposed by [Jiao et al., 2025], leverage a combination of labeled and unlabeled data to enhance model robustness against drift while minimizing the reliance on extensive labeled datasets.

\section{Future Directions and Open Problems}
\label{sec:future_directions}

\#\# Future Directions

As federated learning (FL) continues to evolve in the context of concept drift, several future research directions emerge that address existing gaps in understanding and managing drift in distributed non-stationary environments. This section outlines key areas for further exploration, including the development of a unified theoretical framework, standardized benchmarks, the integration of hybrid adaptation strategies, and the enhancement of model fairness considerations.

\#\#\# Development of a Unified Theoretical Framework

A critical gap in the current literature is the absence of a comprehensive theoretical framework that integrates various methodologies for managing concept drift in federated learning. Most existing studies rely heavily on empirical evaluations, often overlooking the need for a cohesive theoretical model that can unify the diverse approaches to drift management [Mahdi et al., 2025]. Future research should focus on developing theoretical models that account for staggered temporal and spatial heterogeneity in client data distributions. Such frameworks could elucidate how the central server can effectively identify clients experiencing drift and tailor aggregation strategies accordingly [Jothimurugesan et al., 2022]. Additionally, exploring the intersection of model fairness and distributed concept drift presents an opportunity to create frameworks that ensure equitable outcomes across sensitive attributes as data distributions evolve [Salazar et al., 2024].

\#\#\# Standardized Benchmarks for Evaluating Drift-Aware Learning

A significant gap in the current literature is the absence of standardized benchmarks for evaluating federated learning algorithms in the presence of concept drift. Many existing studies rely on fragmented datasets, often utilizing simplistic synthetic adaptations or domain-specific datasets, which do not adequately reflect real-world scenarios [Mahdi et al., 2025]. For instance, while studies have employed datasets like MNIST or CIFAR-10, these do not capture the complexity of natural concept drift that occurs in practical applications, such as IoT environments [Yang et al., 2023]. Future research should prioritize establishing comprehensive benchmark suites that incorporate diverse, real-world datasets exhibiting natural concept drift. This would facilitate rigorous comparative analyses of drift-aware algorithms and contribute to the development of robust FL frameworks.

\#\#\# Integration of Hybrid Adaptation Strategies

Current approaches to managing concept drift in federated learning can be broadly categorized into passive and active/reactive mechanisms. Passive strategies, such as dynamic learning rates, adjust model parameters based on observed performance metrics, while active methods involve explicit detection of drift events [Rahman et al., 2025][Mavromatis et al., 2024]. However, there is a notable lack of studies that effectively combine these strategies to create a more balanced approach to drift adaptation. For example, while passive adaptation mechanisms may be computationally less intensive, they could be less responsive in rapidly changing environments compared to active methods [Mavromatis et al., 2024]. Developing hybrid frameworks that leverage the strengths of both passive and active strategies could enhance the robustness of federated learning systems against concept drift, especially in resource-constrained environments such as mobile and IoT devices [Zhou et al., 2024]. This integration could also address the limitations of existing studies that often assume ideal client capabilities, which may not be realistic in practice.

\#\#\# Addressing Model Fairness in Federated Learning Under Concept Drift

The implications of model fairness in the context of distributed concept drift remain largely unexplored. Future research should investigate how evolving data distributions may lead to biased outcomes, particularly in sensitive applications such as healthcare [CITE:YourPaperID]. To provide a more comprehensive understanding, it is essential to explore strategies that ensure fairness and equity across diverse populations [CITE:YourPaperID]. The main contributions of such research could significantly advance the field of federated learning in the face of concept drift.

\section{Conclusion}
\label{sec:conclusion}

\#\# Conclusion

This survey has provided a comprehensive overview of federated learning (FL) methodologies that address concept drift within distributed non-stationary environments. The analysis highlighted a variety of approaches, including adaptive mechanisms, dynamic learning rates, and the integration of drift detection algorithms, all aimed at enhancing model robustness in the face of evolving data distributions. Notably, frameworks such as FedDrift and FedDrift-Eager exemplify the potential of multi-model strategies to accommodate diverse client data distributions and concept drift scenarios [Jothimurugesan et al., 2022]. Furthermore, innovations like Adaptive-FedAVG and FLAME demonstrate the efficacy of passive and active adaptation techniques in responding to drift, thereby optimizing model accuracy and minimizing degradation over time [Canonaco et al., 2021][Mavromatis et al., 2024].

The contributions of this survey lie in its systematic categorization of existing literature, elucidating the interplay between federated learning and concept drift. By synthesizing findings from over 100 relevant studies, this paper establishes a foundational understanding of the current landscape of drift-aware federated learning, identifying key trends and methodologies that inform future research. Importantly, this survey underscores the necessity of addressing practical challenges such as resource constraints and the demand for real-time adaptability in federated systems, particularly in IoT and edge environments [Benameur et al., 2025][Zhou et al., 2024].

Despite its thoroughness, this survey has limitations that warrant attention. A significant gap exists in the availability of standardized benchmarks for evaluating federated learning algorithms under conditions of concept drift. Most studies rely on proprietary or overly simplistic datasets, which hampers rigorous comparative analyses across different approaches [CITE:high]. Moreover, many methodologies make unrealistic assumptions regarding client capabilities, such as continuous model evaluation or sufficient computational resources, which may not hold true in real-world applications [CITE:medium]. Furthermore, the survey lacks a unified theoretical framework that comprehensively integrates the various methodologies discussed, which could provide deeper insights into the principles guiding federated learning under concept drift [CITE:medium].

\#\#\# Future Research Directions

1. **Development of Standardized Benchmarks**: Establishing unified benchmarks that incorporate natural concept drift scenarios within federated learning contexts is imperative. This would facilitate more rigorous evaluations and comparisons across different algorithms, fostering advancements in the field [CITE:high].

2. **Exploration of Hybrid Adaptation Strategies**: Future research should focus on developing hybrid frameworks that effectively combine client-side and server-side adaptation techniques. Such approaches could leverage the strengths of both passive and active mechanisms, resulting in enhanced reactivity and efficiency in addressing concept drift [CITE:medium].

3. **Theoretical Frameworks for Federated Learning**: There is a pressing need for comprehensive theoretical frameworks that analyze federated learning under distributed concept drift. These frameworks should address staggered temporal and spatial heterogeneity, providing a deeper understanding of the dynamics involved in model aggregation and client drift detection [CITE:medium].

4. **Fairness in Federated Learning**: Investigating the intersection of model fairness and distributed concept drift represents an underexplored area. Research should aim to develop frameworks that ensure equitable outcomes across sensitive attributes as data distributions evolve, thereby addressing potential biases inherent in federated learning systems [CITE:medium].

5. **Scalability of Multi-Model Approaches**: Addressing scalability issues in multi-model approaches remains a critical area for future exploration.

\bibliographystyle{IEEEtran}
\bibliography{references}

\end{document}